%%%%%%%%%%%%%%%%%%%%%%%%%%%%%%%%%%%%%%%%%%%%%%%%%%%%%%%%%%%%%%%%%%%%%%%%%%%%%%%%%%%%%%%%%%%%%%%
%% Credit: Initial version written by [Ivan C. Christov](http://christov.tmnt-lab.org),      %%
%% Purdue University.                                                                        %%
%% License: GPLv3 (see https://www.gnu.org/licenses/gpl-3.0.en.html)                         %%
%%%%%%%%%%%%%%%%%%%%%%%%%%%%%%%%%%%%%%%%%%%%%%%%%%%%%%%%%%%%%%%%%%%%%%%%%%%%%%%%%%%%%%%%%%%%%%%
\DocumentMetadata
 {
 % testphase = phase-I, % tagging without paragraph tagging
 testphase  = phase-II, % tagging with paragraph tagging
 % testphase  = phase-III, % tagging with paragraph sec, toc, blocks and more
 pdfversion = 2.0, % pdfversion must be set here.
 pdfstandard=ua-2, % pdfstandard can be set too
}
\documentclass{article}
\usepackage{amsmath,color}
\usepackage[margin=0.75in]{geometry}

\usepackage{fancyhdr}
\usepackage{totpages}
\fancyhead[L]{Purdue University, ME 50900}
\fancyhead[R]{Prof.\ Ivan C.\ Christov}
\fancyfoot[C]{Page \thepage\ of \ref*{TotPages}}

\DeclareMathOperator{\tr}{tr}

\usepackage[compact]{titlesec}
% redefine (sub)sections to include a period
\renewcommand{\thesubsection}{\arabic{subsection}.}
\renewcommand{\thesubsubsection}{\arabic{subsection}.\arabic{subsubsection}.}

\begin{document}

\thispagestyle{fancy}

\section*{(Some) Vector \& tensor algebra operations}

\subsection{Transpose}
\[
	\text{If}\quad \underline{\underline{T}} = T_{ij}
	\quad\text{then}\quad
	T_{ji} = \underline{\underline{T}}^t \;\; \left[\text{understood as } (\underline{\underline{T}}^t)_{ij \text{ entry}} \right]
\]
\textbf{Remember:} flipping the indices creates a \emph{new} tensor (``the transpose''); $T_{ij} = T_{ji}$ \emph{only} for \emph{symmetric} tensors!\\
(\textbf{Note:} Panton denotes transpose with $^t$, but you might also see $^\top$.)

\subsection{Multiplication}

\subsubsection{vector/vector (dot product; row vector $\times$ column vector = scalar) -- \textit{only one that commutes}}
\[
	\underline{u}\cdot\underline{v} = u_i v_i = v_i u_i = \underline{v}\cdot\underline{u}
\]
\textbf{Note:} For $\underline{v}\cdot\underline{v}=v_iv_i$ \textbf{don't} write $v_i^2$ because this expression has a free index, and it's therefore a vector!

\subsubsection{vector/tensor (row vector $\times$ matrix = row vector, or matrix $\times$ column vector = column vector) -- \textit{does not commute}}
(\textbf{Note:} Panton puts a $\cdot$ between vectors and tensors for clarity; optional when using 1 vs.\ 2 underlines.)
\[
	\underline{v}\,\underline{\underline{T}} = v_i T_{ij}
	\qquad\text{or}\qquad
    (\underline{v}\,\underline{\underline{T}})_{j \text{th entry}} = \sum_{{\color{magenta}i}=1}^3 T_{{\color{magenta}i}j}v_{\color{magenta}i}
\]
but we're allowed to (carefully!) re-index as
\[
    (\underline{v}\,\underline{\underline{T}})_{i \text{th entry}} = \sum_{{\color{magenta}j}=1}^3 T_{{\color{magenta}j}i}v_{\color{magenta}j}
\]
which \textbf{does not commute} because
\[
    \underline{\underline{T}}\,\underline{v} = v_j T_{ij}
	\qquad\text{or}\qquad
    (\underline{\underline{T}}\,\underline{v})_{i \text{th entry}} = \sum_{j=1}^3 T_{ij}v_j
\]
\textbf{Note:} If we always write $T_{ij}$, then you can remember what index to put on $v$ by whether $\underline{v}$ is in-front of $\underline{\underline{T}}$ (then its $v_i$, \textcolor{red}{and $\underline{v}$ is a row vector}) or behind $\underline{\underline{T}}$ (then its $v_j$, \textcolor{red}{and $\underline{v}$ is a column vector}) in the symbolic notation expression.

The following useful fact follows:  $\underline{v}\,\underline{\underline{T}}$ is the row vector with the same entries as the column vector $\underline{\underline{T}}^t\,\underline{v}$.

\subsubsection{vector/vector (dyadic product; column vector $\times$ row vector = matrix) -- \textit{does not commute}}
\[
	\underline{u}\,\underline{v} = u_i v_j
\]
\textbf{does not commute} because
\[
	\underline{v}\,\underline{u} = v_i u_j = u_j v_i
\]
But, if we let $\underline{\underline{T}} =  \underline{u}\,\underline{v}$, then we have shown that $ \underline{v}\,\underline{u} = \underline{\underline{T}}^t$.\\
Sometimes we denote ``$\underline{u}\,\underline{v}$'' as ``$\underline{u}\otimes\underline{v}$'' (so we don't accidentaly confuse it for ``$\underline{u}\cdot\underline{v}$'').

\subsubsection{tensor/tensor (direct product; matrix $\times$ matrix = matrix) -- \textit{does not commute}}
\[
	\underline{\underline{R}} = \underline{\underline{T}}\,\underline{\underline{S}} \quad\text{or}\quad R_{ij} = T_{ik}S_{kj} \quad\Rightarrow\quad\;\text{(for example)}\quad T_{ik}S_{jk} = \underline{\underline{T}}\,\underline{\underline{S}}^t \ne \underline{\underline{R}} \; !
\]
This is simply multiplication of matrices, where $k$ runs over the $i$th row of $\underline{\underline{T}}$ \emph{and} $j$th column of $\underline{\underline{S}}$:
\[
	R_{ij} = \sum_{k=1}^3 T_{ik} S_{kj}
\]
\textcolor{red}{\textbf{``Row times column'' $\Rightarrow$ inner indices} (2nd one on $T$ and 1st one on $S$, in index notation) \textbf{must match!}}

\subsubsection{tensor/tensor (double contraction; ``dot product'' for tensors) -- \textit{does not commute}}
\[
	\underline{\underline{T}} : \underline{\underline{S}} = T_{ij} S_{ji} \quad\Rightarrow\quad\;\text{(for example)}\quad T_{ij} S_{ij} = \underline{\underline{T}} : \underline{\underline{S}}^t \ne \underline{\underline{T}} : \underline{\underline{S}}
\]
\textbf{Note:} If we define $R_{kl} = T_{kj} S_{jl}$ then $\tr(\underline{\underline{R}}) = R_{kk} = R_{ii} = T_{ij} S_{ji}$, in other words $ \underline{\underline{T}} :  \underline{\underline{S}} = \tr( \underline{\underline{T}}\, \underline{\underline{S}})$.

\end{document}
